\scanpage{44}
1 \uline{Rappels} \deleted{\unsure{définitions}} et notations. Opérateurs de Fredholm

2. \uline{\unsure{Produits tensoriels hilbertiens}. \deleted{\unsure{Opérateurs Fredholm}} et de Hilbert-Schmidt}
\margin{\unsure{Traitement} \illegible{3} $\unsure{\Psi}$ est \unsure{vectoriel} \unsure{symétrique}.}

Soit $E$ un espace de \unsure{Banach} (sur le corps des réels ou complexes), désignons par $\overline E\deleted{\unsure{{}^*}}$ l'espace \deleted{\unsure{de base}} vectoriel obtenu en prenant sur $E$ la structure vectorielle \uline{ayant} la même structure additive que $E$, mais où la multiplication scalaire est donnée par $(\lambda,x)\longrightarrow\overline\lambda.x$.
\deleted{\unsure{$\overline E$ est ainsi}}
On a $\lambda\overline x=\overline{\overline\lambda x}$.
On a bien là une \deleted{\unsure{loi de groupe}} \inserted{structure d'espace} vectoriel, \deleted{et la fonction $\|x\|$ sur $\overline E$ est encore une norme}.
L'application canonique \inserted{$x\longrightarrow\overline x$} de $E$ sur $\overline E$ est antilinéaire, i.e \unsure{elle est} additive, et satisfait à :
\[
\overline{\lambda x}=\overline\lambda\,\overline x,
\]
\deleted{\unsure{c'est de plus une isométrie (au point de vue de l'application}\ \unsure{de l'espace de Banach} $\overline E$).}
Notons qu'on a aussi
\[
\overline{\overline E}=E.
\]
\margin{(1) Bien entendu, si le corps des scalaires est le corps des réels, \deleted{$\overline E=E$} $\overline E$ \unsure{possède} la même \unsure{structure} que $E$, \illegible{4}.}

\deleted{$\unsure{E=\overline E}$ \illegible{4}}
\deleted{Si $F$ est un espace vectoriel, une application antilinéaire de $E$ dans $F$ \unsure{est encore une} application linéaire de $\overline E$ dans $F$, \unsure{et de plus} \illegible{2}.}

$E,F$, \deleted{\unsure{étant}} \inserted{\unsure{deux}} espaces vectoriels, une application $(x,y)\longrightarrow u(x,y)$ \deleted{\uline{bilinéaire}} de $E\times F$ \deleted{\unsure{est dite}} dans $\mathbf C$ se dit \uline{sesquilinéaire} si elle est linéaire en $x$, et antilinéaire en $y$ ; cela signifie donc aussi bien qu'elle est une application \uline{bilinéaire} de $E\times\overline F$ dans $\mathbf C$. Ces applications correspondent donc exactement aux applications linéaires de $E\otimes\overline F$ dans $\mathbf C$.
\margin{\unsure{$L^h$}}
\margin{Notons aussi $L(E,F)=L(\overline E,\overline F)$}

\deleted{Si $E$ est un espace de Banach, on \illegible{3} de l'espace $\overline E$ \illegible{4} une norme $\|x\|$ \illegible{2}.}
\mspage{45}
